%!TEX root = ./../../../springer-mono-format/main.tex
%!TEX root = ./../../thesis-trier.tex

\label{def:Sigmaij}
Here, we investigate constraint PDFAs $\automat{B}{P}{\Sigma}{\mu}{1}{F}$
where the state set is $P = \{1,2,\ldots,n\}$ for $n \in \{2,3\}$.
Let
$\Sigma_{i,j} = \{ x \in \Sigma \mid \applytrans{\mu}{i}{x} = j \}$
w.r.t.~$\mathcal B$ for $i,j \in \{1,2,\ldots,n\}$.
Note that, by determinism of $\mathcal B$,
we have $\Sigma_{i,j} \cap \Sigma_{i,k} = \emptyset$
for $j \ne k$ where $i,j,k \in \{1,\ldots,n\}$.

First, in~\cref{constr_sync:sec:two-state}, we investigate the case $P = \{1,2\}$ for an arbitrary alphabet
and then, in~\cref{sec:constr_sync:three_states}, the case $P = \{1,2,3\}$ for $\Sigma= \{a,b\}$.

By choice of notation, $1$ is the start state.
We assume that~$\mathcal B$ is trim.
The case that $\mathcal B$ is strongly connected (including the case $P = \{1\}$, see \cref{constr_sync:cor:one-state})
was already handled in \SEC~\ref{sec:constr_sync:SCC} for general automata.
So, we also assume that $\mathcal B$ is not strongly connected
in the following. 

\subsection{Two-State Constraint PDFAs}
\label{constr_sync:sec:two-state}

\begin{theorem}
	\label{thm:two-state-classification}
	Suppose $P = \{1,2\}$, $\mathcal B$ is trim with start state $1$
	and not strongly connected. Then we have $\Sigma_{1,2} \ne \emptyset$
	and $\Sigma_{2,1} = \emptyset$.
	In this case, by trimness we have $2 \in F$
	and by \cref{constr_sync:prop:final_states,gen_results:cor:factors_etc_complexity} 
	we can assume $F = \{2\}$
	since~$1$ is coaccessible for this final state set.
	Under these assumptions, $L(\mathcal B) = \Sigma_{1,1}^* \Sigma_{1,2} \Sigma_{2,2}^*$ and $\constrsync{L(\mathcal B)}$
	has the following complexities:
	\[
	\begin{cases}
		\text{\PSPACE-complete} & \text{if } ( \Sigma_{1,1} \cup \Sigma_{1,2} ) \setminus \Sigma_{2,2} \ne \emptyset 
		\mbox{ and } \max\{ |\Sigma_{1,1}|, |\Sigma_{2,2}| \} \ge 2, \\
		\text{\NL-complete}        & \text{if } \Sigma_{1,1} \cup \Sigma_{1,2} \subseteq \Sigma_{2,2} \text{ and } |\Sigma_{2,2}|\ge 2, \\ %\max\{ |\Sigma_{1,1}|, |\Sigma_{2,2}| \} \ge 2, \\
		\text{\LOGSPACE-complete}  & \text{if } \max\{ |\Sigma_{1,1}|, |\Sigma_{2,2}| \} = 1, \\
		\text{in $\AC^0$} & \text{if } \Sigma_{1,1} = \Sigma_{2,2} = \emptyset.	
	\end{cases}	
	\]
\end{theorem}
\begin{proof}	
	Under the assumptions from the statement, clearly 
	we have $L(\mathcal B) = \Sigma_{1,1}^* \Sigma_{1,2} \Sigma_{2,2}^*$.
	First, suppose
	\[
	(\Sigma_{1,1} \cup \Sigma_{1,2}) \setminus \Sigma_{2,2} \ne \emptyset
	\text{ and }
	\max\{|\Sigma_{1,1}|, |\Sigma_{2,2}|\} \ge 2.
	\]
	Assume $|\Sigma_{1,1}| \ge 2$.
	Then we can pick distinct $a,b \in \Sigma_{1,1}$
    and $c \in \Sigma_{1,2}$ such that $|\{a,b,c\} \cap \Sigma_{2,2}| \le 2$.  
    We have $\pi_{\{a,b,c\}}(\Sigma_{1,1}^* \Sigma_{1,2} \Sigma_{2,2}^*) = (a+b)^*c\Gamma^*$
    for some $\Gamma \subseteq \{a,b,c\}$ with $|\Gamma| \le 2$.
    By \crefitem{constr_sync:thm:introductory_examples}{constrsync:ex:item:abcL}
    and \cref{thm:hom_lower_bound_complexity} 
    we find that \constrsync{L(\mathcal B)}
    is \PSPACE-hard.
    Next, suppose $|\Sigma_{1,1}| \le 1$ (so that $|\Sigma_{2,2}| \ge 2$).   
    Set 
    \[
    L = \begin{cases}
    	\Sigma_{1,2}\Sigma_{2,2}^* & \text{if } \Sigma_{1,2} \subseteq \Sigma_{2,2}, \\
      	\Sigma_{2,2}^* & \text{otherwise}.
    \end{cases}
    \]
    By assumption, we can pick $a \in \Sigma$
    such that $a \in \Sigma_{1,1} \setminus \Sigma_{2,2}$ if $\Sigma_{1,2} \subseteq \Sigma_{2,2}$
    or $a \in \Sigma_{1,2} \setminus \Sigma_{2,2}$ otherwise.
    For each $v \in \pref{\Sigma_{2,2}} \setminus \Sigma_{2,2} = \{\varepsilon\}$
    the word $va$ is weakly auxiliary state avoiding 
    w.r.t.~$\Sigma_{2,2}$.
    We can pick two distinct $v_1, v_2 \in L$
    and we have $\{v_1, v_2\}^* \subseteq L \subseteq \Sigma_{2,2}^*$.
    Note that $a \in \Sigma_{1,1}$
    implies $\Sigma_{1,1} = \{a\}$ by assumption.
    By \cref{constr_sync:thm:uC}, we find
    that $\constrsync{a^*aL}$
    is \PSPACE-hard for $\Sigma_{1,1} = \{a\}$
    and that $\constrsync{aL}$ is \PSPACE-hard otherwise.
    In the latter case,
    we have $\Sigma_{1,1}^* \Sigma_{1,2} \Sigma_{2,2}^* \cap \Sigma^* a \Sigma^* = \Sigma_{1,1}^* a \Sigma_{2,2}^*$
    where $L = \Sigma_{2,2}^*$. 
    In this case, since $|\Sigma_{1,1}| \le 1$
    and $a \notin \Sigma_{1,1}$
    the problem $\constrsync{\Sigma_{1,1}^* a L}$
    is \PSPACE-hard by \crefitem{constrsync:cor:UanL}{constrsync:cor:UanL:2}
    and so $\constrsync{L(\mathcal B)}$
    is \PSPACE-hard by \cref{constr_sync:gen_results:thm:ideal_separation}.
%        
%   Alter Beweis (auch okay). Neu gemacht am 1.03, so dass weniger andere Sätze benutzt werden.
%    
%    If there exists $a \in \Sigma_{1,2} \setminus \Sigma_{2,2}$,
%    then let $b,c \in \Sigma_{2,2}$ be distinct.
%    As $a \notin \Sigma_{1,1}$,
%    we have $\pi_{\{a,b,c\}}(L(\mathcal B)) \cap \{a,b,c\}^* a  \{a,b,c\}^*
%    = \pi_{\{a,b,c\}}(\Sigma_{1,1})^* a \pi_{\{a,b,c\}}(\Sigma_{2,2})^*$
%    with $\pi_{\{a,b,c\}}(\Sigma_{2,2})^* = \{b,c\}^*$
%    and $\pi_{\{a,b,c\}}(\Sigma_{1,1}) \subseteq \{x\}$ for some $x \in \{b,c\}$.
%    By \cref{constrsync:ex:PSPACEhard}, the problem
%    $\constrsync{ \{\pi_{\{a,b,c\}}(\Sigma_{1,1})^*\} a \{b,c\}^* }$
%    is \PSPACE-hard.
%    So, by \cref{thm:hom_lower_bound_complexity},
%    we find that \constrsync{L(\mathcal B)}
%    is \PSPACE-hard.
%    Next, suppose we have $\Sigma_{1,1} = \{a\}$
%    for some $a \notin \Sigma_{2,2}$ and $\Sigma_{1,2} \subseteq \Sigma_{2,2}$.
%    Let $L = \Sigma_{1,2}\Sigma_{2,2}^*$ and $v_1, v_2 \in L$ be distinct.
%    Then $wa$ is weakly auxiliary state avoiding w.r.t.~$\Sigma_{2,2}$
%    for each $w \in \pref{\Sigma_{2,2}} \setminus \Sigma_{2,2}$.
%    As $v_1, v_2 \in \Sigma_{2,2}^*$, we have $\{v_1, v_2\}^* \subseteq L \subseteq \pref{\Sigma_{2,2}^*}$.
%    By \cref{constr_sync:thm:uC}, we find
%    that \constrsync{a^*aL}
%    is \PSPACE-hard.
    Containment in \PSPACE\  follows from~\cref{constr_sync:thm:PSPACE}.    
    
    Next, suppose $\Sigma_{1,1} \cup \Sigma_{1,2} \subseteq \Sigma_{2,2}$
    with $|\Sigma_{2,2}| \ge 2$.
    Then 
    $\fact{L(\mathcal B)} = \fact{\Sigma_{2,2}^*}$.
    By \cref{gen_results:cor:factors_etc_complexity}, 
    $\constrsync{L(\mathcal B)} \equiv_m^{\AC^0} \constrsync{\Sigma_{2,2}^*}$,
    and the latter problem is \NL-complete 
    by \cref{constr_sync:cor:one-state}. 
    Now, suppose $\max\{ |\Sigma_{1,1}|, |\Sigma_{2,2}| \} = 1$.
    Then, by \cref{constr_sync:thm:infinite}  and 
    \cref{constr_sync:thm:sparse_basic_L_cases},
    the problem is \LOGSPACE-complete.
    Finally, for $\Sigma_{1,1} = \Sigma_{2,2} = \emptyset$,
    $L(\mathcal B)$ is finite
    and containment in $\AC^0$ follows from \cref{constr_sync:thm:finite}. 	
\end{proof}


\begin{corollary}
	\label{constr_sync:cor:two-state_binary}
For a two-state constraint automaton $\mathcal B$
over a binary alphabet, \constrsync{L(\mathcal B)}
is in \NL.
\end{corollary}
\begin{proof}
	We can assume that $\mathcal B$ is trim with start state $1$. Then we have $\Sigma_{1,2} \ne \emptyset$.
	So, as $\Sigma_{1,1} \cap \Sigma_{1,2} = \emptyset$,
	we have  $|\Sigma_{1,1}|\le 1$. Then, for $|\Sigma_{2,2}| \le 1$,
	by \cref{thm:two-state-classification} 
	the problem is in \LOGSPACE. 
	For $|\Sigma_{2,2}| = 2$,
	we have $\Sigma_{1,1} \cup \Sigma_{1,2}\subseteq \Sigma_{2,2}$
	and the problem is in \NL\ 
	by \cref{thm:two-state-classification}.
\end{proof}

\subsection{Three-State Constraint PDFAs Over a Binary Alphabet}
\label{sec:constr_sync:three_states}

Here, we prove the complexity classification
for $P = \{1,2,3\}$ with $\Sigma = \{a,b\}$.
Broadly, we structure the proof by distinguishing the different cases for the strongly connected components of $\mathcal B$.

\begin{theorem} 
	\label{constr_sync:thm:three-states}
	Suppose $P = \{1,2,3\}$, $\mathcal B$ is not strongly connected and $\Sigma = \{a,b\}$.
	Let $\mathcal B$ be trim with start state $1$. Then
	we have $\Sigma_{1,2} \cup \Sigma_{1,3} \ne \emptyset$ (and so $|\Sigma_{1,1}| \le 1$ as $\mathcal B$ is deterministic). 
	%Without loss of generality, assume $a \in \Sigma_{1,2}$.
	If $L(\mathcal B)$ is finite, then $\constrsync{L(\mathcal B)}$
	is in $\AC^0$ (\THM~\ref{constr_sync:thm:finite}). So, we assume $L(\mathcal B)$
	is infinite in the following.
	
	\begin{enumerate}		
		\item Suppose $\{1\}, \{2\}, \{3\}$ are the strongly connected components. Then we must have $\Sigma_{2,3} = \emptyset$
		or $\Sigma_{3,2} = \emptyset$.
		Without loss of generality, assume $\Sigma_{3,2} = \emptyset$.
		Then the following complexities arise (recall that $\mathcal B$ is trim):
		
		\begin{center} 
			\begin{tabular}{l|l}
				$\max\{ |\Sigma_{2,2}|, |\Sigma_{3,3}| \} = 2$ & \NL-complete \\ \hline 
				$\max\{ |\Sigma_{2,2}|, |\Sigma_{3,3}| \} \le 1$ & \multirow{ 2}{*}{\NP-complete} \\
				$\quad\land |\Sigma_{2,2}| = 1 \land \Sigma_{2,3} \ne \emptyset \land (\Sigma_{1,1} \cup \Sigma_{1,2}) \setminus \Sigma_{2,2} \ne \emptyset$ & \\ \hline 
				$\max\{ |\Sigma_{2,2}|, |\Sigma_{3,3}| \} \le 1$ & \multirow{ 2}{*}{\LOGSPACE-complete} \\
				$\quad\land ( \Sigma_{2,2} = \emptyset \lor \Sigma_{2,3} = \emptyset \lor (\Sigma_{1,1} \cup \Sigma_{1,2}) \subseteq \Sigma_{2,2} )$ & \\ \hline 
				%\multicolumn{2}{c}{\itshape \small Complexities in case $\{1\}, \{2\}$, $\{3\}$ are the strongly connected components in $\mathcal B$.}
			\end{tabular}
        \end{center}		
		
		\item Suppose $\{1,2\}$ is a strongly connected component, or equivalently
		$\Sigma_{1,2}$ and $\Sigma_{2,1}$ are both non-empty. 
		Then the following complexities arise:	
				
		\begin{center} 
		\begin{tabular}{l|l}
			$\Sigma_{3,3} = \{a,b\}$ & \NL-complete \\ \hline
			$|\Sigma_{3,3}| \le 1$,
			$\Sigma_{1,1} = \Sigma_{2,2} = \emptyset$,
			$|\Sigma_{1,2}| = |\Sigma_{2,1}| = 1$ & \LOGSPACE-complete \\ \hline
			All other cases (for the stated assumptions) & \PSPACE-complete
		\end{tabular}
		\end{center}		
		
		\item Suppose  $\{2,3\}$ is a strongly connected component, or equivalently
		 $\Sigma_{2,3}$ and $\Sigma_{3,2}$ are both non-empty.
		 Then the following complexities arise:		 
	
		\begin{tabular}{l|p{6cm}}
			$|\Sigma_{2,3}| + |\Sigma_{2,2}| + |\Sigma_{3,2}| + |\Sigma_{3,3}| = 2$ & \LOGSPACE-complete \\ \hline 
			$|\Sigma_{2,3}| + |\Sigma_{2,2}| + |\Sigma_{3,2}| + |\Sigma_{3,3}| = 4$ & \NL-complete \\ \hline 
			$|\Sigma_{2,3}| + |\Sigma_{2,2}| + |\Sigma_{3,2}| + |\Sigma_{3,3}| = 3$ & All possible constraint automata, up to symmetry, and their complexities
			are listed in Table~\ref{constr_sync:table:three_states_23} and~\ref{tab:binary_automata_three_states}.		
		\end{tabular} 
		
	\end{enumerate}
\end{theorem}
\begin{proof} For $\max\{ |\Sigma_{2,2}|, |\Sigma_{3,3}| \} = 2$, $\mathcal B$ has a trap state.
	Then, by~\crefitem{constrsync:cor:upperset}{constrsync:cor:upperset:item:trimtrap},
	we get that $\constrsync{L(\mathcal B)}$
	is \NL-complete.
	So, for the remaining proof,
	we assume $\max\{ |\Sigma_{2,2}|, |\Sigma_{3,3}| \} \le 1$.
	Also, we  assume that $L(\mathcal B)$ is infinite,
	which gives \LOGSPACE-hardness by~\cref{constr_sync:thm:infinite},
	since the cases when~$L(\mathcal B)$ is finite are easy to
	identify, given by trim PDFAs without any loops.	
	
\medskip 
\noindent\textbf{\texorpdfstring{$\{1\}$, $\{2\}$, $\{3\}$}{\{1\},\{2\},\{3\}} are the Strongly Connected Components.}
Without loss of generality, suppose $\Sigma_{3,2} = \emptyset$.
As $\mathcal B$ 
is trim, the states $2$ and $3$ are reachable from~$1$, 
i.e., $\Sigma_{1,2} \ne \emptyset \land ( \Sigma_{1,3} \ne \emptyset \lor \Sigma_{2,3} \ne \emptyset)$, which
implies $|\Sigma_{1,1}|\le 1$.
Then, by \cref{constr_sync:move_states}
we can assume $F = \{2,3\}$, so that
\begin{equation}\label{constrsync:eqn:threestatesSCC}
	L(\mathcal B) = \Sigma_{1,1}^* \Sigma_{1,2} \Sigma_{2,2}^* \Sigma_{2,3} \Sigma_{3,3}^* \cup \Sigma_{1,1}^* \Sigma_{1,2} \Sigma_{2,2}^*
	\cup \Sigma_{1,1}^* \Sigma_{1,3} \Sigma_{3,3}^*. 
\end{equation}
Hence
$L(\mathcal B) \subseteq \Sigma_{1,1}^*\Sigma_{1,2}^* \Sigma_{1,3}^* \Sigma_{2,2}^* \Sigma_{2,3}^* \Sigma_{3,3}^*$.
As $\max\{ |\Sigma_{2,2}|, |\Sigma_{3,3}| \} \le 1$
the language~$L(\mathcal B)$ is letter-bounded.
Now, suppose $( \Sigma_{1,1} \cup \Sigma_{1,2} ) \setminus \Sigma_{2,2} \ne \emptyset$
and there exist $b,c \in \Sigma$ such that 
$b \in \Sigma_{2,2}$ and $c \in \Sigma_{2,3}$.
Then $ub^*c \cap L(\mathcal B)$ is infinite for some
$u \in \Sigma_{1,1}\Sigma_{1,2}$ with $u \notin b^*$.
By \cref{constr_sync:thm:letter_bounded} 
the problem \constrsync{L(\mathcal B)} is \NP-complete.
Otherwise, we have $\fact{L(\mathcal B)} = \fact{L}$
where (for non-mutually exclusive cases)
\[
L = \begin{cases}
	\Sigma_{2,2}^* \Sigma_{2,3} \Sigma_{3,3}^* \cup \Sigma_{2,2}^* \cup \Sigma_{1,1}^* \Sigma_{1,3} \Sigma_{3,3}^*
	& \text{if } \Sigma_{1,1} \cup \Sigma_{1,2} \subseteq \Sigma_{2,2}, \\ % , \Sigma_{1,2} \ne \emptyset, \\
	%\Sigma_{1,1}^* \cup \Sigma_{1,1}^* \Sigma_{1,3} \Sigma_{3,3}^*
	%& \text{if } \Sigma_{1,2} = \emptyset, \\
	\Sigma_{1,1}^* \Sigma_{1,2} \Sigma_{2,3} \Sigma_{3,3}^* \cup \Sigma_{1,1}^* \Sigma_{1,2}\cup \Sigma_{1,1}^* \Sigma_{1,3} \Sigma_{3,3}^* & \text{if } \Sigma_{2,2} = \emptyset, \\
	\Sigma_{1,1}^* \Sigma_{1,2} \Sigma_{2,2}^* \cup \Sigma_{1,1}^* \Sigma_{1,2}\cup \Sigma_{1,1}^* \Sigma_{1,3} \Sigma_{3,3}^* &
	\text{if } \Sigma_{2,3} = \emptyset.
\end{cases}
\]
As $\max\{ |\Sigma_{1,1}|, |\Sigma_{2,2}|, |\Sigma_{3,3}| \} \le 1$, for the listed cases
$L$ equals a finite union of languages 
of the form $u^* v w^*$ for $u,v,w \in \Sigma^*$.
By \COR~\ref{gen_results:cor:factors_etc_complexity}, \cref{constr_sync:thm:union,constr_sync:thm:sparse_basic_L_cases}
it follows that \constrsync{L(\mathcal B)} is in \LOGSPACE.   

\medskip
\noindent\textbf{\texorpdfstring{$\{1,2\}$}{{1,2}} is a Strongly Connected Component.} 
In this case, since $\{1,2\}$ contains the start state and~$\mathcal B$ is trim,
we must have $3 \in F$ and that $3$ is reachable from every other state.
Thus, by \crefitem{constr_sync:move_states}{constr_sync:move_states:2}
we can assume $F = \{3\}$. 
Without loss of generality, suppose $\Sigma_{1,3} = \{a\}$
and $\Sigma_{1,2} = \{b\}$, as for $\Sigma_{1,3} = \emptyset$
it is $\Sigma_{2,3} \ne \emptyset$ 
and, by \cref{constr_sync:prop:final_states}, we can consider $2$ as
the start state. Then $\Sigma_{1,1} = \emptyset$
and
\[
L(\mathcal B) = 
\begin{cases}
	(by)^*(a + b\Sigma_{2,3})\Sigma_{3,3}^*  & \text{if } \Sigma_{2,1} = \{y\} \text{ and } \Sigma_{2,2} = \emptyset, \\
	(bb+ba)^*a\Sigma_{3,3}^*  & \text{if } \Sigma_{2,1} = \{a,b\}, \\
	(ba^*b)^*a\Sigma_{3,3}^*  & \text{if } \Sigma_{2,2} = \{a\} \text{ and } \Sigma_{2,1} = \{b\}, \\
	(bb^*a)^*a\Sigma_{3,3}^*    & \text{if } \Sigma_{2,2} = \{b\} \text{ and } \Sigma_{2,1} = \{a\}. 
\end{cases}
\]
As $|\Sigma_{3,3}| \le 1$ by assumption,
with \cref{constr_sync:cor:u_star_finite_v_star} we find that \constrsync{L(\mathcal B)} is in \LOGSPACE.
For the remaining cases, as $\fact{(ba^*b)a\Sigma_{3,3}^*} = \fact{(a+bb)ba\Sigma_{3,3}^*}$
and $\fact{(bb^*a)a\Sigma_{3,3}^*} = \fact{(b+ab)aa\Sigma_{3,3}^*}$,
\PSPACE-completeness follows from~\cref{constrsync:prop:PSPACE-hard-cases}, \ref{constrsync:Cu:1},\ref{constrsync:Cu:2} \& \ref{constrsync:Cu:3}.

\medskip
\noindent\textbf{\texorpdfstring{$\{2,3\}$}{{2,3}} is a Strongly Connected Component.} 
Since~$\mathcal B$ is trim with start state~$1$, we have $F \cap \{2,3\} \ne \emptyset$.
Thus, by \cref{constr_sync:move_states:1,constr_sync:move_states:2} of \cref{constr_sync:move_states} 
we can suppose $F = \{3\}$.
Now, first, let $|\Sigma_{2,3}| + |\Sigma_{2,2}| + |\Sigma_{3,2}| + |\Sigma_{3,3}| = 2$.
Then $\Sigma_{1,1}=\Sigma_{2,2}=\emptyset$
and $\Sigma_{2,3}=\{x\}, \Sigma_{3,2}=\{y\}$ for some $x,y \in \Sigma$,
which gives
$L(\mathcal B) = \Sigma_{1,1}^* \Sigma_{1,2} (xy)^* \cup \Sigma_{1,1}^*\Sigma_{1,3} (yx)^* = \Sigma_{1,1}^* ( \Sigma_{1,2} \cup \Sigma_{1,3} ) (yx)^*$.
As $|\Sigma_{1,1}| \le 1$, by \cref{constr_sync:cor:u_star_finite_v_star},
\constrsync{L(\mathcal B)} is in \LOGSPACE.
Since the $L(\mathcal B)$ is infinite,
the problem is \LOGSPACE-hard 
by \cref{constr_sync:thm:infinite}.
Next, suppose $|\Sigma_{2,3}| + |\Sigma_{2,2}| + |\Sigma_{3,2}| + |\Sigma_{3,3}| = 4$ and set $S = \{2,3\}$.
Then $\mu(p, x)$ is defined with $\mu(p, x) \in S$ for each $p \in S$
and $x \in \Sigma$.
As $\mathcal B$ is trim, by~\crefitem{constrsync:cor:upperset}{constrsync:cor:upperset:item:S},
\constrsync{L(\mathcal B)} is \NL-complete.
Lastly, suppose 
$|\Sigma_{2,3}| + |\Sigma_{2,2}| = |\Sigma_{3,2}| + |\Sigma_{3,3}| = 3$, or equivalently
that precisely three transitions are defined between the states in $\{2,3\}$.
Then, all possible combinations and the resulting complexities
where, without loss of generality, we assume $a \in \Sigma_{1,2}$
are listed in \cref{constr_sync:table:three_states_23,tab:binary_automata_three_states}. 
\Cref{constr_sync:table:three_states_23} 
gives the proofs. 

%!TEX root = ../../../springer-mono-format/book.tex
%!TEX root = ../../thesis-trier.tex

\newcommand{\textXabebXeae}{
	$L_{1}^{\{2\}} \subseteq L_2^{\{2\}}$
}
\newcommand{\textXabebXbae}{
	$L_{1}^{\{2\}} \subseteq \suff{L_2^{\{2\}}}$
}
\newcommand{\textXabebXeabe}{
	$L_{1}^{\{2\}} \subseteq \suff{L_2^{\{2\}}}$
}
\newcommand{\textXabebXeab}{
	$a \in \suff{L_2^{\{2,3\}}}$, 
	$b \in \suff{L_{3}^{\{2,3\}}}$
}

\newcommand{\textXabeaXeae}{
	$L_{1}^{\{2\}} \subseteq L_2^{\{2\}}$
}
\newcommand{\textXabeaXbae}{
	$\Sigma^*bba\Sigma^* \cap L_1^{\{2\}} = b^*bba(a+ba)^*$
}
\newcommand{\textXabeaXeabe}{
	$\Sigma^*bba\Sigma^* \cap 
	L_1^{\{2\}}
	= bba(a+ba)^*$
}
\newcommand{\textXabeaXeab}{
	$a \in \suff{L_2^{\{2,3\}}}$,
	$b \in \suff{L_3^{\{2,3\}}}$
}

\newcommand{\textXbaebXeae}{
	$\Sigma^*aab\Sigma^* \cap L_{1}^{\{2\}} = aab(b+ab)^*$
}
\newcommand{\textXbaebXbae}{
	$\Sigma^*aab\Sigma^* \cap L_{1}^{\{2\}} = b^*aab(b+ab)^*$
}
\newcommand{\textXbaebXeabe}{
	$\Sigma^*aab\Sigma^* \cap L_{1}^{\{2\}} = aab(b+ab)^*$
}
\newcommand{\textXbaebXeab}{
	$\Sigma^*aab\Sigma^* \cap L_{1}^{\{2\}} = aab(b+ab)^*$
}

\newcommand{\textXbaeaXeae}{
	$L_{1}^{\{2\}} \subseteq \suff{L_2^{\{2\}}}$
}
\newcommand{\textXbaeaXbae}{
	$\Sigma^*bab\Sigma^* \cap L_{1}^{\{2\}} = b^*bab(b+aa)^*$
}
\newcommand{\textXbaeaXeabe}{
	$L_{1}^{\{2\}} \subseteq \suff{L_2^{\{2\}}}$
}
\newcommand{\textXbaeaXeab}{
	$\Sigma^*bab\Sigma^* \cap L_{1}^{\{2\}} = bab(b+aa)^*$
}

\newcommand{\textXeabeaXeae}{
	$L_{1}^{\{2\}} \subseteq \suff{L_2^{\{2\}}}$
}
\newcommand{\textXeabeaXbae}{
	$\Sigma^*bba\Sigma^* \cap L_{1}^{\{2\}} = b^*bba(aa+ba)^*$
}
\newcommand{\textXeabeaXeabe}{
	$\Sigma^*bba\Sigma^* \cap L_{1}^{\{2\}} = bba(aa+ba)^*$
}
\newcommand{\textXeabeaXeab}{
	$a \in \suff{L_{2}^{\{2,3\}}}$,
	$b \in \suff{L_{3}^{\{2,3\}}}$
}

\newcommand{\textXeabebXeae}{
	$\Sigma^*aab \Sigma^* \cap L_{1}^{\{2\}} = aab(bb+ab)^*$
}
\newcommand{\textXeabebXbae}{
	$\Sigma^*aab \Sigma^* \cap L_{1}^{\{2\}} = b^*aab(bb+ab)^*$
}
\newcommand{\textXeabebXeabe}{
	$\Sigma^*aab \Sigma^* \cap L_{1}^{\{2\}} = aab(bb+ab)^*$
}
\newcommand{\textXeabebXeab}{
	$\Sigma^*aab \Sigma^* \cap L_{1}^{\{2\}} = aab(bb+ab)^*$
}

\newcommand{\textXeaabXeae}{
	$L_{1}^{\{3\}} \subseteq \suff{L_3^{\{3\}}}$
}
\newcommand{\textXeaabXbae}{
	$\Sigma^*bba\Sigma^*\cap L_{1}^{\{2\}} = b^*bba(aa^*b)^*$
}
\newcommand{\textXeaabXeabe}{
	$L_1^{\{3\}} \subseteq \suff{L_3^{\{3\}}}$
}
\newcommand{\textXeaabXeab}{
	$\Sigma^*bba\Sigma^*\cap L_{1}^{\{2\}}) = bba(aa^*b)^*$
}

\newcommand{\textXeabaXeae}{
	$L_1^{\{2\}} \subseteq \suff{L_2^{\{2\}}}$
}
\newcommand{\textXeabaXbae}{
	$L_1^{\{2\}} \subseteq \suff{L_2^{\{2\}}}$
}
\newcommand{\textXeabaXeabe}{
	$\Sigma^*bab\Sigma^* \cap L_{1}^{\{3\}} = bab(b+aa)^*$
}
\newcommand{\textXeabaXeab}{
	$L_1^{\{3\}} \subseteq \suff{L_3^{\{3\}}}$
}

\newcommand{\textXebabXeae}{
	$\Sigma^*aba\Sigma^* \cap L_{1}^{\{3\}} = aba(a+bb)^*$
}
\newcommand{\textXebabXbae}{
	$\Sigma^*aba\Sigma^* \cap L_{1}^{\{3\}} = b^*aba(a+bb)^*$
}
\newcommand{\textXebabXeabe}{
	$\Sigma^*aba\Sigma^* \cap L_{1}^{\{3\}} = aba(a+bb)^*$
}
\newcommand{\textXebabXeab}{
	$\Sigma^*aba\Sigma^* \cap L_{1}^{\{3\}} = aba(a+bb)^*$
}

\newcommand{\textXebbaXeae}{
	$L_1^{\{2\}} \subseteq \suff{L_2^{\{2\}}}$
}
\newcommand{\textXebbaXbae}{
	$L_1^{\{2\}} \subseteq \suff{L_2^{\{2\}}}$
}
\newcommand{\textXebbaXeabe}{
	$L_1^{\{3\}} \subseteq \suff{L_3^{\{3\}}}$
}
\newcommand{\textXebbaXeab}{
	$L_1^{\{3\}} \subseteq \suff{L_3^{\{3\}}}$
}

\newcommand{\textXeaeabXeae}{
	$L_{1}^{\{2\}} \subseteq \suff{L_2^{\{2\}}}$
}
\newcommand{\textXeaeabXbae}{
	$\Sigma^* baab \Sigma^* \cap L_{1}^{\{2\}} = b^*baab(aa+ab)^*$
}
\newcommand{\textXeaeabXeabe}{
	$L_{1}^{\{2\}} \subseteq \suff{L_2^{\{2\}}}$
}
\newcommand{\textXeaeabXeab}{
	$\Sigma^* bba \Sigma^* \cap L_{1}^{\{3\}} = bba(aa+ba)^*$
}

\newcommand{\textXebeabXeae}{
	$L_{1}^{\{2\}} \subseteq \suff{L_2^{\{2\}}}$
}
\newcommand{\textXebeabXbae}{
	$L_{1}^{\{2\}} \subseteq \suff{L_2^{\{2\}}}$
}
\newcommand{\textXebeabXeabe}{
	$L_{1}^{\{2\}} \subseteq \suff{L_2^{\{2\}}}$
}
\newcommand{\textXebeabXeab}{
	$L_{1}^{\{3\}} \subseteq \suff{L_3^{\{3\}}}$
}

 \label{def:LiFnotation}
\begin{longtable}{l | l | l | l}
	\caption[3-state binary constraint automaton complexities, case $\{2,3\}$]{Set $L_i^F = L(\mathcal B_{i,F})$. 
		Complexities for $|\Sigma_{2,2}| + |\Sigma_{2,3}| + |\Sigma_{3,3}| + |\Sigma_{3,2}| = 3$ (see Thm.~\ref{constr_sync:thm:three-states}). 
		\PSPACE-completeness follows from 
		Thm.~\ref{constr_sync:gen_results:thm:ideal_separation}, Prop.~\ref{constrsync:prop:PSPACE-hard-cases}
		and Cor.~\ref{constrsync:cor:UanL}
		and membership in \NL from Cor.~\ref{constr_sync:cor:UVW} and Thm.~\ref{constr_sync:thm:union}
		(the arguments to apply these statements are indicated in the third column).
		Up to symmetry, we assume $a \in \Sigma_{1,2}$. 
		By giving the subalphabets $\Sigma_{i,j}$, 
		the first column lists all (up to symmetry) automata~$\mathcal B_{|\{2,3\}}$,		 
		and the second column lists all possibilities to enter the strongly connected component $\{2,3\}$
		from~$1$.
		The combinations give (up to symmetry) all possibilities  (see also Table~\ref{tab:binary_automata_three_states}).
	} \label{constr_sync:table:three_states_23} \\
	\toprule
	$\Sigma_{2,2}, \Sigma_{2,3}, \Sigma_{3,3}, \Sigma_{3,2}$ & $\Sigma_{1,1}, \Sigma_{1,2}, \Sigma_{1,3}$ & Argument & Complexity \\
	\hline
	\endfirsthead
	
	\toprule
	$\Sigma_{2,2}, \Sigma_{2,3}, \Sigma_{3,3}, \Sigma_{3,2}$ & $\Sigma_{1,1}, \Sigma_{1,2}, \Sigma_{1,3}$ & Argument & Complexity \\
	\hline
	\endhead
	
	\hline
	\multicolumn{4}{r}{Continued on next page} \\
	\endfoot
	
	\bottomrule
	\endlastfoot
	
	$\{a\},\{b\},\emptyset,\{b\}$ &
	$\emptyset,\{a\},\emptyset$ & {\textXabebXeae} & \NL \\
	& $\{b\},\{a\},\emptyset$ & {\textXabebXbae} & \NL \\
	& $\emptyset,\{a,b\},\emptyset$ & {\textXabebXeabe} & \NL \\
	& $\emptyset,\{a\},\{b\}$ & {\textXabebXeab} & \NL \\
	
	\hline
	
	$\{a\},\{b\},\emptyset,\{a\}$ &
	$\emptyset,\{a\},\emptyset$ & {\textXabeaXeae} & \NL \\
	& $\{b\},\{a\},\emptyset$ & {\textXabeaXbae} & \PSPACE-c \\
	& $\emptyset,\{a,b\},\emptyset$ & {\textXabeaXeabe} & \PSPACE-c \\
	& $\emptyset,\{a\},\{b\}$ & {\textXabeaXeab} & \NL \\
	
	\hline
	
	$\{b\},\{a\},\emptyset,\{b\}$ &
	$\emptyset,\{a\},\emptyset$ & {\textXbaebXeae} & \PSPACE-c \\
	& $\{b\},\{a\},\emptyset$ & {\textXbaebXbae} & \PSPACE-c \\
	& $\emptyset,\{a,b\},\emptyset$ & {\textXbaebXeabe} & \PSPACE-c \\
	& $\emptyset,\{a\},\{b\}$ & {\textXbaebXeab} & \PSPACE-c \\
	
	\hline
	
	$\{b\},\{a\},\emptyset,\{a\}$ &
	$\emptyset,\{a\},\emptyset$ & {\textXbaeaXeae} & \NL \\
	& $\{b\},\{a\},\emptyset$ & {\textXbaeaXbae} & \PSPACE-c \\
	& $\emptyset,\{a,b\},\emptyset$ & {\textXbaeaXeabe} & \NL \\
	& $\emptyset,\{a\},\{b\}$ & {\textXbaeaXeab} & \PSPACE-c \\
	
	\hline
	
	$\emptyset,\{a,b\},\emptyset,\{a\}$ &
	$\emptyset,\{a\},\emptyset$ & {\textXeabeaXeae} & \NL \\
	& $\{b\},\{a\},\emptyset$ & {\textXeabeaXbae} & \PSPACE-c \\
	& $\emptyset,\{a,b\},\emptyset$ & {\textXeabeaXeabe} & \PSPACE-c \\
	& $\emptyset,\{a\},\{b\}$ & {\textXeabeaXeab} & \NL \\
	
	\hline
	
	$\emptyset,\{a,b\},\emptyset,\{b\}$ &
	$\emptyset,\{a\},\emptyset$ & {\textXeabebXeae} & \PSPACE-c \\
	& $\{b\},\{a\},\emptyset$ & {\textXeabebXbae} & \PSPACE-c \\
	& $\emptyset,\{a,b\},\emptyset$ & {\textXeabebXeabe} & \PSPACE-c \\
	& $\emptyset,\{a\},\{b\}$ & {\textXeabebXeab} & \PSPACE-c \\
	
	\hline
	
	$\emptyset,\{a\},\{a\},\{b\}$ &
	$\emptyset,\{a\},\emptyset$ & {\textXeaabXeae} & \NL \\
	& $\{b\},\{a\},\emptyset$ & {\textXeaabXbae} & \PSPACE-c \\
	& $\emptyset,\{a,b\},\emptyset$ & {\textXeaabXeabe} & \NL \\
	& $\emptyset,\{a\},\{b\}$ & {\textXeaabXeab} & \PSPACE-c \\
	
	\hline
	
	$\emptyset,\{a\},\{b\},\{a\}$ &
	$\emptyset,\{a\},\emptyset$ & {\textXeabaXeae} & \NL \\
	& $\{b\},\{a\},\emptyset$ & {\textXeabaXbae} & \NL \\
	& $\emptyset,\{a,b\},\emptyset$ & {\textXeabaXeabe} & \PSPACE-c \\
	& $\emptyset,\{a\},\{b\}$ & {\textXeabaXeab} & \NL \\
	
	\hline
	
	$\emptyset,\{b\},\{a\},\{b\}$ &
	$\emptyset,\{a\},\emptyset$ & {\textXebabXeae} & \PSPACE-c \\
	& $\{b\},\{a\},\emptyset$ & {\textXebabXbae} & \PSPACE-c \\
	& $\emptyset,\{a,b\},\emptyset$ & {\textXebabXeabe} & \PSPACE-c \\
	& $\emptyset,\{a\},\{b\}$ & {\textXebabXeab} & \PSPACE-c \\
	
	\hline
	
	$\emptyset,\{b\},\{b\},\{a\}$ &
	$\emptyset,\{a\},\emptyset$ & {\textXebbaXeae} & \NL \\
	& $\{b\},\{a\},\emptyset$ & {\textXebbaXbae} & \NL \\
	& $\emptyset,\{a,b\},\emptyset$ & {\textXebbaXeabe} & \NL \\
	& $\emptyset,\{a\},\{b\}$ & {\textXebbaXeab} & \NL \\
	
	\hline
	
	$\emptyset,\{a\},\emptyset,\{a,b\}$ &
	$\emptyset,\{a\},\emptyset$ & {\textXeaeabXeae} & \NL \\
	& $\{b\},\{a\},\emptyset$ & {\textXeaeabXbae} & \PSPACE-c \\
	& $\emptyset,\{a,b\},\emptyset$ & {\textXeaeabXeabe} & \NL \\
	& $\emptyset,\{a\},\{b\}$ & {\textXeaeabXeab} & \PSPACE-c \\
	
	\hline
	
	$\emptyset,\{b\},\emptyset,\{a,b\}$ &
	$\emptyset,\{a\},\emptyset$ & {\textXebeabXeae} & \NL \\
	& $\{b\},\{a\},\emptyset$ & {\textXebeabXbae} & \NL \\
	& $\emptyset,\{a,b\},\emptyset$ & {\textXebeabXeabe} & \NL \\
	& $\emptyset,\{a\},\{b\}$ & {\textXebeabXeab} & \NL \\
\end{longtable}
%\newcommand{\clg}{\cellcolor{green!15}}
%\newcommand{\cdg}{\cellcolor{green!35}}
%\newcommand{\clr}{\cellcolor{red!15}}
%\newcommand{\cdr}{\cellcolor{red!35}}
\newcommand{\cdr}[1]{#1}
\newcommand{\clr}[1]{#1}
\newcommand{\cdg}[1]{#1}
\newcommand{\clg}[1]{#1}


%
%\NewDocumentCommand{\fillpattern}{O{north west lines} O{blue!50} m m}{%
%	\draw[pattern=#1, pattern color=#2] (#3.north west)rectangle (#4.south east);
%}


% 
% Teilalphabete \Sigma_ii als Argumente; 0 bedeutet Teilalphabet ist leer
\newcommand{\drawaut}[8]{
\scalebox{.76}{
	\begin{tikzpicture}[>=latex',shorten >=1pt,node distance=1.5cm,on grid,auto,baseline=0,initial text={}]
	\tikzset{every state/.style={minimum size=1pt,fill=white}}
	
	%\tikzset{background grid/.style={thick,draw=red,step=.5cm}	show background grid}

	\ifthenelse { \equal {#8} {dr} } { \draw[pattern=north east lines, pattern color=red!45]   (-1.8,-0.9) rectangle (1.8,1);  } { 	}
    \ifthenelse { \equal {#8} {lr} } { \draw[pattern=horizontal lines, pattern color=red!35]   (-1.8,-0.9) rectangle (1.8,1); } {  }
    \ifthenelse { \equal {#8} {dg} } { \draw[pattern=north west lines, pattern color=green!45] (-1.8,-0.9) rectangle (1.8,1);} { 	 }
    \ifthenelse { \equal {#8} {lg} } { \draw[pattern=vertical lines,   pattern color=green!35] (-1.8,-0.9) rectangle (1.8,1); } {  }		
	
	\node[state] (1) {};
	\node[state] (2) [right of=1] {};
	%\node[state,initial,initial where=below] (0) [left of=1] {};
	\node[state] (0) [left of=1] {};
			
  	\ifthenelse { \equal {#1} {0} } {} { \path[->] (0) edge [loop above, above] node {$#1$} (0); }
  	\ifthenelse { \equal {#2} {0} } {} { \path[->] (1) edge [loop above] node {$#2$} (1); }
 	\ifthenelse { \equal {#3} {0} } {} { \path[->] (2) edge [loop above, above] node {$#3$} (2); }
  	\ifthenelse { \equal {#4} {0} } {} { \path[->] (0) edge [bend right,below] node {$#4$} (2); }
	    
  	\path[->] (0) edge node {$#5$} (1);
  	\path[->] (1) edge [bend left]  node {$#6$} (2);
  	\path[->] (2) edge [bend left]  node {$#7$} (1);
  
\end{tikzpicture}}
}



% https://latexdraw.com/how-to-create-a-lined-paper-background-in-latex-using-tikz/
% https://tex.my/2020/05/24/creating-tiled-background-patterns/
% https://tex-stackexchange-com.translate.goog/questions/67092/set-table-background-hatched-and-shaded-using-tikz?_x_tr_sl=en&_x_tr_tl=de&_x_tr_hl=de&_x_tr_pto=sc
% https://tex.stackexchange.com/questions/352229/pattern-background-in-table-cell-or-minipage
% https://tex.stackexchange.com/questions/67092/set-table-background-hatched-and-shaded-using-tikz
% https://tex.stackexchange.com/questions/193018/scan-a-cell-of-a-table-with-lines
% https://tex.stackexchange.com/questions/502474/create-a-table-with-once-and-twice-hatched-regions-and-curly-braces
% https://tex.stackexchange.com/questions/545899/how-to-hatch-pattern-this-table-excluding-the-circled-region

\begin{table}[ht]
    \centering
     \caption[3-state binary constraint PDFA pictures, case \{2,3\}]{\footnotesize 
    	$3$-state PDFAs with at least one undefined transition among the states, numbered $1,2,3$ from left to right, in the strongly connected component $\{2,3\}$.
    	%where the states are numbered $1,2,3$ from left to right
    	The cases are the same as in Table~\ref{constr_sync:table:three_states_23}.
    	% By \cref{constr_sync:move_states}
    	% any non-empty subset of $\{2,3\}$ as final states gives the same complexity;
    	% so final states are not marked.
    	The start state is~$1$.
    	Final states are not marked, as by \cref{constr_sync:move_states}
    	any non-empty subset of $\{2,3\}$
    	gives the same complexity. 
    	%is indicated by an incoming arrow.
    }
   % vielleicht besser so
%    \begin{tikzpicture}
%    	\matrix (m)[matrix of nodes, style={nodes={rectangle,draw,minimum width=3em}}, minimum height=3em, row sep=-\pgflinewidth, column sep=-\pgflinewidth,ampersand replacement =\&]
%    	{
%    		1 \& 2 \& 3 \\
%    		4 \& 5 \& 6 \\
%    		7 \& 8 \& 9 \\
%    	};
%    	
%    	\begin{pgfonlayer}{myback}
%    		\fillpattern{m-3-2}{m-3-2}
%    		\fillpattern[dots][green!50!black]{m-3-1}{m-3-1}
%    		\fillpattern[sixpointed stars][violet!50]{m-3-3}{m-3-3}
%    	\end{pgfonlayer}
%    \end{tikzpicture}

\begin{tabular}{cccc} %\hline
\drawaut{0}{a}{0}{0}{a}{b}{b}{lg}    & \drawaut{b}{a}{0}{0}{a}{b}{b}{lg}   & \drawaut{0}{a}{0}{0}{a,b}{b}{b}{lg} & \drawaut{0}{a}{0}{b}{a}{b}{b}{dg} \\ 
\drawaut{0}{a}{0}{0}{a}{b}{a}{lg}    & \drawaut{b}{a}{0}{0}{a}{b}{a}{dr}   & \drawaut{0}{a}{0}{0}{a,b}{b}{a}{lr}  & \drawaut{0}{a}{0}{b}{a}{b}{a}{dg} \\ 
\drawaut{0}{b}{0}{0}{a}{a}{b}{lr}    & \cdr{\drawaut{b}{b}{0}{0}{a}{a}{b}{dr}}   & \clr{\drawaut{0}{b}{0}{0}{a,b}{a}{b}{lr}}  & \clr{\drawaut{0}{b}{0}{b}{a}{a}{b}{lr}} \\ 
\clg{\drawaut{0}{b}{0}{0}{a}{a}{a}{lg}}    & \cdr{\drawaut{b}{b}{0}{0}{a}{a}{a}{dr}}   & \clg{\drawaut{0}{b}{0}{0}{a,b}{a}{a}{lg}}  & \clr{\drawaut{0}{b}{0}{b}{a}{a}{a}{lr}} \\ \hline
\drawaut{0}{0}{0}{0}{a}{a,b}{a}{lg} & \drawaut{b}{0}{0}{0}{a}{a,b}{a}{dr}  & \drawaut{0}{0}{0}{0}{a,b}{a,b}{a}{lr} & \drawaut{0}{0}{0}{b}{a}{a,b}{a}{dg} \\ 
\drawaut{0}{0}{0}{0}{a}{a,b}{b}{lr} & \drawaut{b}{0}{0}{0}{a}{a,b}{b}{dr}  & \drawaut{0}{0}{0}{0}{a,b}{a,b}{b}{lr} & \drawaut{0}{0}{0}{b}{a}{a,b}{b}{lr} \\ 
\clg\drawaut{0}{0}{a}{0}{a}{a}{b}{lg}     & \cdr\drawaut{b}{0}{a}{0}{a}{a}{b}{dr}    & \clg\drawaut{0}{0}{a}{0}{a,b}{a}{b}{lg}   & \cdr\drawaut{0}{0}{a}{b}{a}{a}{b}{dr}   \\ 
\clg\drawaut{0}{0}{b}{0}{a}{a}{a}{lg}     & \clg\drawaut{b}{0}{b}{0}{a}{a}{a}{lg}  & \clr\drawaut{0}{0}{b}{0}{a,b}{a}{a}{lr}   & \clg\drawaut{0}{0}{b}{b}{a}{a}{a}{lg}   \\ \hline
\clr  \drawaut{0}{0}{a}{0}{a}{b}{b}{lr}   & \cdr\drawaut{b}{0}{a}{0}{a}{b}{b}{dr}    & \clr\drawaut{0}{0}{a}{0}{a,b}{b}{b}{lr}  & \clr\drawaut{0}{0}{a}{b}{a}{b}{b}{lr}   \\ 
\clg\drawaut{0}{0}{b}{0}{a}{b}{a}{lg}     & \clg\drawaut{b}{0}{b}{0}{a}{b}{a}{lg}  & \clg\drawaut{0}{0}{b}{0}{a,b}{b}{a}{lg}   & \clg\drawaut{0}{0}{b}{b}{a}{b}{a}{lg}   \\ 
    \clg\drawaut{0}{0}{0}{0}{a}{a}{a,b}{lg}   & \cdr\drawaut{b}{0}{0}{0}{a}{a}{a,b}{dr}& \clg\drawaut{0}{0}{0}{0}{a,b}{a}{a,b}{lg} & \clr\drawaut{0}{0}{0}{b}{a}{a}{a,b}{lr} \\ 
    \clg\drawaut{0}{0}{0}{0}{a}{b}{a,b}{lg}   & \drawaut{b}{0}{0}{0}{a}{b}{a,b}{lg}  & \clg\drawaut{0}{0}{0}{0}{a,b}{b}{a,b}{lg} & \clg\drawaut{0}{0}{0}{b}{a}{b}{a,b}{lg} \\ \hline
    \end{tabular}
   
 \label{tab:binary_automata_three_states}
\end{table}

\Cref{tab:binary_automata_three_states}
gives graphical depictions of the resulting constraint automata
for the cases from \cref{constr_sync:table:three_states_23}.
\end{proof}

With \cref{thm:two-state-classification}, we need at least three states
to realize \NP-complete constrained synchronization problems.
Over a binary alphabet, 
by \cref{constr_sync:cor:two-state_binary,constr_sync:thm:three-states},
three states are the least number of states to realize \PSPACE-complete
constrained synchronization problems. 

